\documentclass[10pt,a4paper]{article}
\usepackage{xeCJK}
\setCJKmainfont{Yu Mincho}
\setCJKsansfont{Yu Gothic}
\setCJKmonofont{Yu Gothic}
\usepackage[margin=25mm]{geometry}
\usepackage{amsmath,amssymb,amsthm,bm}
\usepackage{enumerate}
\usepackage{graphicx}
\usepackage{tikz}
\usepackage{csquotes}
\usepackage[backend=biber,style=apa,sorting=nyt]{biblatex}
\addbibresource{manuscript/references.bib}
\usetikzlibrary{arrows.meta,positioning,calc,fit,shapes.geometric}

\tikzset{
  dagobs/.style={draw, rounded corners=1pt, minimum width=10mm, minimum height=7mm, align=center, inner sep=2pt, font=\small},
  daglat/.style={draw, circle, minimum size=8mm, align=center, inner sep=1pt, font=\small},
  dagerr/.style={draw, circle, dashed, minimum size=7mm, align=center, inner sep=1pt, font=\scriptsize},
  dagmiss/.style={draw, diamond, aspect=1.7, minimum width=10mm, inner sep=1pt, align=center, font=\small},
  dagcond/.style={dagobs, double, double distance=1pt},
  dagarrow/.style={-{Latex[length=2.2mm,width=1.5mm]}, line width=.55pt},
  dagdashed/.style={dagarrow, dashed},
  dagnone/.style={line width=.45pt, dotted},
  dagnote/.style={align=center, font=\scriptsize, inner sep=1pt}
}

\newtheorem{assumption}{仮定}
\newtheorem{theorem}{定理}
\newtheorem{proposition}{命題}
\newtheorem{lemma}{補題}
\newtheorem{corollary}{系}
\newtheorem{remark}{注意}
\newtheorem{definition}{定義}

\newcommand{\indep}{\mathrel{\perp\!\!\!\perp}}
\newcommand{\E}{\mathbb{E}}
\newcommand{\Prb}{\mathbb{P}}
\newcommand{\R}{\mathbb{R}}
\newcommand{\1}{\mathbf{1}}
\newcommand{\calO}{\mathcal{O}}
\newcommand{\calF}{\mathcal{F}}
\newcommand{\calV}{\mathcal{V}}
\newcommand{\calY}{\mathcal{Y}}
\newcommand{\Y}{\bm{Y}}
\newcommand{\D}{\bm{D}}
\newcommand{\eps}{\bm{\varepsilon}}

\title{欠測IVと潜在変数モデリングによる非無作為欠測の識別}
\author{%
本多将大\\[-1mm]
\small 慶應義塾大学大学院経済学研究科
\and
星野崇宏\\[-1mm]
\small 慶應義塾大学経済学部
\and
Taisuke Otsu\\[-1mm]
\small Department of Economics, London School of Economics\\[-1mm]
\small Houghton Street, London WC2A 2AE, UK\\[-1mm]
\small \texttt{t.otsu@lse.ac.uk}\quad \texttt{http://econ.lse.ac.uk/staff/otsu\_taisuke/}
}
\date{}

\begin{document}
\maketitle

\begin{abstract}
非無作為欠測（nonignorable missingness; NI/MNAR）では，欠測確率が未観測値に依存するため，通常の多重代入や観測変数のみに基づくIPWでは母集団分布や平均の識別が困難である。本研究では，d'Haultfoeuille (2010) の shadow variable / missing IV 識別を，スカラーの欠測アウトカムからベクトル値の欠測変数 $\Y=(Y_1,\ldots,Y_K)'$ へ拡張する。第一段階では，欠測指標 $R$ が $\Y$ に依存しても，$\Y$ を条件づけると常時観測補助変数 $V$ が欠測機構から排除される $R\indep V\mid \Y$ と vector completeness の下で，欠測確率 $\pi(\Y)=\Prb(R=1\mid \Y)$，同時分布 $P(\Y,V)$，および平均ベクトル $\E(\Y)$ が識別されることを示す。第二段階では，復元された $P(\Y,V)$ に $\Y=\bm h(F)+\eps,\ F=g(V)+U$ という潜在測定モデルの識別条件を課すことで，$g$ と $h_1,\ldots,h_K$ の識別を議論する。したがって本稿の主識別結果は，潜在変数に依存するNMARを直接モデル化する方法ではなく，d'Haultfoeuille 型の欠測IV識別のベクトル拡張である。
\end{abstract}

\section{Introduction}

非無作為欠測（non-MAR, NI/MNAR）は，観測データからの識別を困難にする代表的問題である。Rubin 型の欠測分類では，欠測がデータに依存しない場合を MCAR，観測変数のみに依存する場合を MAR，未観測値そのものに依存する場合を NI/MNAR と呼ぶ。MAR では，観測変数に条件づけることで多重代入（MI）や inverse probability weighting（IPW）が利用できる。一方，NI/MNAR では，欠測確率が未観測値に依存するため，観測データだけでは母集団分布や平均は一般に識別されない。

例えば，$R=1$ のときのみ $Y$ が観測される場合，完全ケース分布は
\begin{equation}
  p(y\mid R=1)=
  \frac{\Prb(R=1\mid Y=y)p_Y(y)}
       {\int \Prb(R=1\mid Y=t)p_Y(t)dt}
  \label{eq:intro-selected-density}
\end{equation}
であり，母集団密度 $p_Y$ と選択確率 $\Prb(R=1\mid Y)$ は積としてのみ現れる。したがって，欠測機構または母集団分布に追加制約を置かなければ，$p_Y$ も $\E(Y)$ も一意には定まらない。

この非識別性に対する一つの有力な解決策が，常に観測される補助変数を用いた shadow variable アプローチである。本稿の出発点は，d'Haultfoeuille (2010, \textit{Journal of Econometrics}) の ``A new instrumental method for dealing with endogenous selection'' である。同論文の識別戦略は，通常の処置効果IVのように「操作変数がアウトカムから排除される」ことではなく，「アウトカムを条件づけたとき，操作変数が選択・欠測指標から排除される」ことに基づく。d'Haultfoeuille (2010) における $Z$ は，後続研究では shadow variable，nonresponse instrument，または missing-data instrumental variable と呼ばれる変数に対応する。通常の処置効果推定における操作変数とは異なり，ここでの IV は欠測機構から排除された補助変数である。

本稿の主な貢献は，d'Haultfoeuille 型の shadow variable 識別を，スカラーの欠測変数からベクトル値の欠測変数へ拡張し，さらに共通潜在因子構造と結合する点にある。従来の素朴な拡張では，各欠測変数 $Y_j$ に対してそれぞれ別の shadow variable を用意する必要があるように見える。これに対して本稿では，複数の欠測変数が共通潜在因子 $F$ を通じて同時に生成される場合，スカラーまたはベクトル値の常時観測補助変数 $V$ によってベクトル $\Y=(Y_1,\ldots,Y_K)'$ の joint distribution を識別できる条件を示す。

この点は，欠測変数間の相関を明示的に認めるためにも重要である。旧稿では $Y_j$ ごとに
\begin{equation}
  D_j\indep V\mid Y_j
\end{equation}
を置き，さらに同時観測確率を $\prod_j \pi_j(Y_j)$ と表記していた。しかしこの書き方は，欠測機構および $Y_j$ 間の関係を過度に独立化しているように見える。ベクトル $\Y$ に対して議論すれば，$Y_1,\ldots,Y_K$ の相関を許したまま，
\begin{equation}
  R\indep V\mid \Y
  \label{eq:intro-vector-shadow}
\end{equation}
を中心条件として識別論を構成できる。

本稿では，まずIV型構造
\begin{equation}
  F=g(V)+U
  \label{eq:intro-iv-latent}
\end{equation}
に集中する。ここで $V$ は常時観測される景気先行指標，$F$ は不況などの共通潜在要因，$\Y$ は複数企業の株価指数・財務指標・調査指標などと解釈できる。欠測または選択が各 $Y_j$ に発生する場合でも，共通の $F$ を介して $V$ が $\Y$ 全体に情報を持つならば，$V$ はベクトル $\Y$ に対する shadow variable として機能し得る。ただし，ベクトル値 $\Y$ に対する completeness を満たすには，$V$ のサポートと $\Y$ との関連性が十分に豊かである必要がある。

重要な区別として，本稿の主結果は欠測が $\Y$ に依存する
\begin{equation}
  R^*=k(\Y)+\eta
\end{equation}
場合に関する。欠測が潜在因子そのものに依存する $R^*=k(F)+\eta$ の場合には，一般に $R\indep V\mid \Y$ は成立せず，d'Haultfoeuille 型の証明は適用できない。この場合は Appendix \ref{app:latent-factor-missingness} で非識別または別アプローチの問題として扱う。

一方，$V=g(F)+U$ とする proxy 型は，測定誤差モデルや因子分析の文脈では興味深いが，欠測補完の識別戦略としては追加検討を要する。したがって本稿では，proxy 型を中心的主張には含めず，今後の課題として位置づける。ただし，旧稿で検討した proxy 型の補足的識別議論は，Appendix \ref{app:proxy-type} に残す。

本稿の構成は d'Haultfoeuille (2010) の構成を参考にする。第2節で識別問題を定式化し，主定理を示す。第3節で推定方法を述べる。第4節ではシミュレーション設計を提示する。第5節では実証分析または擬似実証の雛形を示す。第6節で結論を述べ，すべての技術的証明は Appendix A にまとめる。追加測定誤差，経験ベイズ的解釈，潜在因子依存欠測，および proxy 型は Appendix B--E にまとめる。

\section{Identification}

\subsection{Existing shadow-variable strategy}

\begin{figure}[tbp]
\centering
\begin{tikzpicture}[x=1cm,y=1cm]
  \node[dagobs] (Z) at (0,0) {$Z$};
  \node[dagerr] (eps) at (1.6,1.25) {$\epsilon$};
  \node[dagobs] (Y) at (3.0,0) {$Y$};
  \node[dagerr] (eta) at (4.5,1.25) {$\eta$};
  \node[dagmiss] (R) at (6.0,0) {$R$};
  \draw[dagarrow] (Z) -- node[dagnote,above] {$\phi$} (Y);
  \draw[dagarrow] (eps) -- (Y);
  \draw[dagarrow] (Y) -- node[dagnote,above] {$\psi$} (R);
  \draw[dagarrow] (eta) -- (R);
  \node[dagnote] at (3,-1.0) {$R=\psi(Y,\eta),\quad Y=\phi(Z,\epsilon),\quad \eta\indep (Z,\epsilon)$};
  \node[dagnote] at (3,-1.55) {排除制約：$Z\nrightarrow R$，したがって $R\indep Z\mid Y$};
  \draw[dagnone] (Z.south) .. controls +(1.8,-1.15) and +(-1.8,-1.15) .. (R.south);
\end{tikzpicture}
\caption{d'Haultfoeuille 型 shadow variable 識別のDAG。$Z$ は $Y$ とは関連するが，$Y$ を条件付けると欠測指標 $R$ から排除される。}
\label{fig:dag-dhaultfoeuille}
\end{figure}

$R=1$ のときに $Y$ が観測される状況を考える。d'Haultfoeuille 型の識別では，常に観測される，または分布が識別可能な変数 $Z$ が
\begin{equation}
  R\indep Z\mid Y
  \label{eq:dh-exclusion}
\end{equation}
を満たし，かつ $Y$ について十分な情報を持つことを利用する。例えば
\begin{equation}
  Y=\phi(Z,\epsilon),\qquad R=\psi(Y,\eta),\qquad \eta\indep (Z,\epsilon)
\end{equation}
ならば，\eqref{eq:dh-exclusion} が成立する。

d'Haultfoeuille (2010) では，$D$ は選択ダミー，$Y$ は欠測し得るアウトカム，$Z$ は操作変数である。観測構造は，$D=1$ のとき $Y$ が観測され，$Z$ の分布は本標本または補助情報から識別されている，というものである。本稿では欠測データ文脈に合わせて，選択ダミー $D$ を完全ケース指標 $R$，操作変数 $Z$ を常時観測補助変数 $V$，スカラーアウトカム $Y$ をベクトル値欠測変数 $\Y$ に置き換える。

欠測確率を
\begin{equation}
  \pi(Y)=\Prb(R=1\mid Y)
\end{equation}
とおくと，shadow variable 条件の下で
\begin{equation}
  \E\left\{\frac{R}{\pi(Y)}\mid Z\right\}=1
  \label{eq:dh-correct-moment}
\end{equation}
が成り立つ。ここで分母は $P(Y)$ ではなく，欠測確率 $\pi(Y)$ である。さらに completeness
\begin{equation}
  \E\{a(Y)\mid Z\}=0 \quad\Rightarrow\quad a(Y)=0
  \label{eq:dh-completeness}
\end{equation}
と positivity $\pi(Y)>0$ の下で，$\pi(Y)$ と $P(Y,Z)$ が識別される。

Zhao and Shao (2015) は，共変量 $X$ と shadow variable $Z$ が観測され，アウトカム $Y$ が非無視可能欠測する GLM を考える。仮定
\begin{equation}
  P(R\mid Y,Z,X)=P(R\mid Y,X)
\end{equation}
の下で，
\begin{equation}
P(Z\mid Y,X,R=1)=P(Z\mid Y,X)
\end{equation}
が成り立つ。さらに
\begin{equation}
P(Z\mid Y,X)=
\frac{P(Y\mid Z,X)P(Z\mid X)}{\int P(Y\mid z,X)P(z\mid X)dz}
\end{equation}
であるため，$P(Y\mid Z,X)$ に GLM 等を仮定すれば，この条件付き分布に基づく擬似尤度を構成できる。

\subsection{Definition and basic equation}

\begin{definition}[Shadow variable]
\label{def:shadow-variable}
$R=1$ のときのみ $Y$ が観測され，共変量 $X$ と補助変数 $Z$ は常に観測されるとする。$Z$ が $Y$ に対する shadow variable であるとは，次を満たすことをいう。
\begin{enumerate}[(i)]
  \item 除外制約：
  \begin{equation}
    R\indep Z\mid (Y,X).
    \label{eq:shadow-exclusion}
  \end{equation}
  \item 関連性：
  \begin{equation}
    Z\not\indep Y\mid X.
    \label{eq:shadow-relevance}
  \end{equation}
\end{enumerate}
さらにノンパラメトリック識別には，関連性を十分に強めた completeness 条件が必要となる。
\end{definition}

重要なのは，shadow variable が MAR を仮定するための変数ではない点である。MAR は
\begin{equation}
  \Prb(R=1\mid Y,X,Z)=\Prb(R=1\mid X,Z)
\end{equation}
を要求するが，shadow variable アプローチでは
\begin{equation}
  \Prb(R=1\mid Y,X,Z)=\Prb(R=1\mid Y,X)
\end{equation}
を許す。すなわち欠測は $Y$ に依存してよい。識別に必要なのは，$Y$ を条件づけた後に $Z$ が欠測に追加的な影響を持たないことと，$Z$ が $Y$ に十分な情報を持つことである。

共変量 $X$ を省略し，$R=1$ のときのみ $Y$ が観測され，$Z$ は常に観測されるとする。$\pi(y)=\Prb(R=1\mid Y=y)$，$w(y)=1/\pi(y)$ とおく。$R\indep Z\mid Y$ の下では
\begin{align}
  \E\{Rw(Y)\mid Z\}
  &=\E\left[\E\{Rw(Y)\mid Y,Z\}\mid Z\right] \\
  &=\E\left[w(Y)\Prb(R=1\mid Y,Z)\mid Z\right] \\
  &=\E\left[w(Y)\pi(Y)\mid Z\right]=1.
  \label{eq:shadow-moment-basic}
\end{align}
したがって，未知関数 $w$ は観測分布から構成される積分方程式
\begin{equation}
  \E\{Rw(Y)\mid Z\}=1
  \label{eq:shadow-integral-equation}
\end{equation}
の解として特徴づけられる。

\begin{assumption}[Shadow completeness]
\label{ass:shadow-completeness-review}
任意の可積分関数 $a$ について，
\begin{equation}
  \E\{a(Y)\mid Z\}=0\quad \mathrm{a.s.}
  \quad\Longrightarrow\quad
  a(Y)=0\quad \mathrm{a.s.}
  \label{eq:shadow-completeness-review}
\end{equation}
が成り立つ。
\end{assumption}

\begin{proposition}[欠測確率の識別]
\label{prop:shadow-propensity-id}
$R\indep Z\mid Y$，positivity $\Prb(R=1\mid Y)>0$，および仮定 \ref{ass:shadow-completeness-review} が成り立つとする。このとき，$w(Y)=1/\Prb(R=1\mid Y)$，したがって $\pi(Y)$ は観測分布から識別される。
\end{proposition}

$\pi$ が識別されれば，任意の可積分関数 $\varphi$ について
\begin{equation}
  \E\{\varphi(Y,Z)\}
  =\E\left\{\frac{R\varphi(Y,Z)}{\pi(Y)}\right\}
  \label{eq:shadow-ipw-functional}
\end{equation}
が成立する。特に $\E(Y)=\E\{RY/\pi(Y)\}$ である。

\subsection{Our vector model}

本節の主対象は，d'Haultfoeuille (2010) の $(D,Y,Z)$ を
\begin{equation}
  D\mapsto R,\qquad
  Y\mapsto \Y,\qquad
  Z\mapsto V
  \label{eq:dh-to-vector-map}
\end{equation}
と置き換えたベクトル版である。ここで $R$ は完全ケース指標，$\Y$ は同時分布を復元したい欠測変数ベクトル，$V$ は常時観測される shadow variable である。潜在因子 $F$ は，$V$ と $\Y$ の関連性および vector completeness を支える構造として導入する。

\begin{figure}[tbp]
\centering
\begin{tikzpicture}[x=1cm,y=1cm]
  \node[dagobs] (V) at (0,0) {$V$};
  \node[dagerr] (U) at (1.5,1.25) {$U$};
  \node[daglat] (F) at (2.0,0) {$F$};
  \node[dagobs] (Y1) at (4.0,1.0) {$Y_1$};
  \node[dagobs] (Y2) at (4.0,0) {$Y_2$};
  \node[dagobs] (Y3) at (4.0,-1.0) {$Y_3$};
  \node[dagmiss] (R) at (6.2,0) {$R$};
  \node[dagerr] (eta) at (6.2,1.25) {$\eta$};
  \draw[dagarrow] (V) -- node[dagnote,above] {$g$} (F);
  \draw[dagarrow] (U) -- (F);
  \draw[dagarrow] (F) -- node[dagnote,above] {$h_1$} (Y1);
  \draw[dagarrow] (F) -- node[dagnote,above] {$h_2$} (Y2);
  \draw[dagarrow] (F) -- node[dagnote,below] {$h_3$} (Y3);
  \draw[dagarrow] (Y1) -- (R);
  \draw[dagarrow] (Y2) -- (R);
  \draw[dagarrow] (Y3) -- (R);
  \draw[dagarrow] (eta) -- (R);
  \node[dagnote] at (3.1,-2.05) {$\Y=(Y_1,Y_2,Y_3)'$ は相関してよい。欠測指標 $R$ は $\Y$ 全体に依存してよい。};
  \node[dagnote] at (3.1,-2.55) {$\Y$ で条件付けると $V\to F\to \Y\to R$ が遮断され，$R\indep V\mid \Y$。};
\end{tikzpicture}
\caption{本稿の基本モデル。常時観測補助変数 $V$ が潜在因子 $F$ を介してベクトル $\Y$ と関連し，欠測指標 $R$ は $\Y$ 全体に依存する。}
\label{fig:dag-vector-basic}
\end{figure}

$K\geq 2$ とし，ベクトル値の測定変数を
\begin{equation}
  \Y=(Y_1,\ldots,Y_K)'\in\R^K
\end{equation}
と定義する。潜在因子 $F$ に対する測定モデルを
\begin{equation}
  Y_j=h_j(F)+\varepsilon_j,
  \qquad j=1,\ldots,K,
  \label{eq:vector-measurement}
\end{equation}
またはベクトル表記で
\begin{equation}
  \Y=\bm h(F)+\eps,
  \qquad \bm h(F)=\{h_1(F),\ldots,h_K(F)\}'
  \label{eq:vector-measurement-compact}
\end{equation}
と書く。識別の正規化として，主に $h_1(f)=f$ を置く。代替的には，$\E(F)=0$，$\operatorname{Var}(F)=1$，かつ $h_1$ が単調増加であるという正規化を用いてもよい。

本稿の中心となるIV型構造は
\begin{equation}
  F=g(V)+U,
  \label{eq:main-iv-structure}
\end{equation}
である。ここで $V$ はスカラーでもベクトルでもよい。後者の場合，$V\in\R^q$ は複数の常時観測された先行指標または補助変数を表す。記法上は単に $V$ と書くが，ベクトル値 $\Y$ に対する completeness を満たすには，$V$ のサポートと $\Y$ との関連性が十分に豊かである必要がある。

各 $Y_j$ は欠測し得るが，以下では joint distribution の識別を明確にするため，まず完全観測パターンを表す指標
\begin{equation}
  R=\1\{Y_1,\ldots,Y_K\text{ がすべて観測される}\}
\end{equation}
を用いる。欠測機構は
\begin{equation}
  R=\1\{R^*>0\},
  \qquad
  R^*=k(\Y)+\eta
  \label{eq:vector-y-missing}
\end{equation}
であり，欠測確率を
\begin{equation}
  \pi(\bm y)=\Prb(R=1\mid \Y=\bm y)
  =\Prb\{\eta>-k(\bm y)\}
  \label{eq:vector-pi}
\end{equation}
と定義する。この定式化では，欠測が $Y_j$ ごとに独立であるとは仮定しない。$Y_j$ 間の相関も，欠測確率が $\Y$ の複数成分に同時に依存することも許す。

観測データは
\begin{equation}
  \calO=\{V,R,R\Y,\D\}
  \label{eq:vector-observed-data}
\end{equation}
である。ここで $\D=(D_1,\ldots,D_K)'$ は項目別の観測指標であり，$R=\prod_{j=1}^K D_j$ である。以下の識別定理では $R=1$ の完全ケースを用いるが，部分観測パターンを利用する拡張は別途，パターン別の欠測確率 $\pi_{\bm d}(\bm y)=\Prb(\D=\bm d\mid \Y=\bm y)$ により定式化できる。

\begin{assumption}[外生性]
\label{ass:vector-exogeneity}
$V,U,\eps,\eta$ は互いに独立であり，$F=g(V)+U$，$\Y=\bm h(F)+\eps$，$R^*=k(\Y)+\eta$ が成り立つ。
\end{assumption}

\begin{assumption}[positivity]
\label{ass:vector-positivity}
ある定数 $c>0$ が存在して
\begin{equation}
  \pi(\Y)\geq c\quad \mathrm{a.s.}
  \label{eq:vector-positivity}
\end{equation}
が成り立つ。
\end{assumption}

\begin{assumption}[vector shadow completeness]
\label{ass:vector-completeness}
任意の可積分関数 $a:\R^K\to\R$ について，
\begin{equation}
  \E\{a(\Y)\mid V\}=0\quad \mathrm{a.s.}
  \quad\Longrightarrow\quad
  a(\Y)=0\quad \mathrm{a.s.}
  \label{eq:vector-completeness}
\end{equation}
が成り立つ。
\end{assumption}

\begin{assumption}[潜在測定モデルの識別条件]
\label{ass:vector-latent-identification}
復元された同時分布 $P(\Y,V)$ に対し，
\begin{equation}
  \Y=\bm h(F)+\eps,
  \qquad
  F=g(V)+U
  \label{eq:vector-latent-id}
\end{equation}
が，正規化 $h_1(f)=f$ の下で一意に識別されると仮定する。これは，非線形測定誤差モデルまたは Hu--Schennach 型の識別定理に対応する仮定である。
この仮定は d'Haultfoeuille (2010) の shadow-variable 識別そのものではなく，第一段階で復元された $P(\Y,V)$ に対して潜在測定関数を回復するための追加的な測定モデル識別条件である。
\end{assumption}

\subsection{Main identification results}

\begin{lemma}[ベクトル条件付き独立性]
\label{lem:vector-shadow-independence}
仮定 \ref{ass:vector-exogeneity} と欠測機構 \eqref{eq:vector-y-missing} の下で，
\begin{equation}
  R\indep V\mid \Y
  \label{eq:vector-shadow-independence}
\end{equation}
が成り立つ。
\end{lemma}

\begin{theorem}[欠測確率と joint distribution の識別]
\label{thm:vector-shadow-identification}
仮定 \ref{ass:vector-exogeneity}--\ref{ass:vector-completeness} の下で，欠測確率 $\pi(\Y)$，同時分布 $P(\Y,V)$，および平均ベクトル $\E(\Y)$ は観測分布から識別される。
\end{theorem}

定理 \ref{thm:vector-shadow-identification} が，本稿における d'Haultfoeuille (2010) の直接のベクトル拡張である。次の定理は，この識別結果で得られた $P(\Y,V)$ を入力として，潜在測定モデルの関数を回復する第二段階の結果である。

\begin{theorem}[IV型潜在測定関数の識別]
\label{thm:vector-iv-functions}
\eqref{eq:vector-measurement-compact}，\eqref{eq:main-iv-structure}，\eqref{eq:vector-y-missing}，および仮定 \ref{ass:vector-exogeneity}--\ref{ass:vector-latent-identification} が成り立つとする。このとき，$\E(\Y)$，$g$，および $h_1,\ldots,h_K$ は観測分布から識別される。
\end{theorem}

\begin{corollary}[二次元の場合]
\label{cor:two-dimensional}
$K=2$ とし，$\Y=(Y_1,Y_2)'$ とする。欠測指標を $R=1\{Y_1,Y_2\text{ がともに観測}\}$ とし，$R^*=k(Y_1,Y_2)+\eta$ を許す。このとき $Y_1$ と $Y_2$ は $F$ を通じて相関し，欠測も両方の値に依存してよい。それでも $R\indep V\mid (Y_1,Y_2)$ と vector completeness が成立すれば，$P(Y_1,Y_2,V)$ は識別される。これは旧稿のような $\pi_1(Y_1)\pi_2(Y_2)$ 型の積重みを必要としない。
\end{corollary}

\begin{remark}[旧稿からの修正点]
旧稿では，各 $Y_j$ について $D_j\indep V\mid Y_j$ を置き，$\prod_j\pi_j(Y_j)$ 型の重みで同時分布を復元していた。しかし，この表記は $Y_j$ 間の相関や欠測指標間の依存を不要に制限しているように見える。本稿では，$\Y$ 全体を一つのベクトルとして扱い，$R\indep V\mid \Y$ と $\pi(\Y)=\Prb(R=1\mid \Y)$ を用いる。これにより，$Y_1,\ldots,Y_K$ の相関を認めたまま，joint distribution $P(\Y,V)$ の識別を示すことができる。
\end{remark}

\section{Estimation}

本節では，識別結果を実装するための推定手順を整理する。d'Haultfoeuille (2010) と同様に，中心となる対象は選択確率の逆数
\begin{equation}
  w(\bm y)=\frac{1}{\pi(\bm y)}
\end{equation}
であり，これは条件付きモーメント
\begin{equation}
  \E\{Rw(\Y)\mid V\}=1
  \label{eq:estimation-main-equation}
\end{equation}
の解として特徴づけられる。推定は大きく，$w$ の第一段階推定，IPW による $P(\Y,V)$ の復元，復元分布に基づく潜在測定関数の推定，の三段階からなる。

\subsection{Parametric or sieve estimation of the selection probability}

第一の実装方法は，$w(\bm y)$ または $\pi(\bm y)$ に柔軟なパラメトリックまたは sieve 表現を置くことである。例えば
\begin{equation}
  w(\bm y;\alpha)=1+\exp\{s(\bm y)'\alpha\}
  \label{eq:sieve-weight}
\end{equation}
と置く。ここで $s(\bm y)$ は多項式，スプライン，基底展開，またはニューラルネットワーク等であってよい。$w(\bm y)\geq 1$ は $0<\pi(\bm y)\leq 1$ に対応する。

$\alpha$ は，任意の補助関数 $m(V)$ に対して
\begin{equation}
  \E\left[\{Rw(\Y;\alpha)-1\}m(V)\right]=0
  \label{eq:gmm-weight}
\end{equation}
を満たすように GMM で推定できる。$m(V)$ と $s(\Y)$ の次元を増やせば，非パラメトリックな近似に近づく。ただし，$\Y$ が高次元の場合には ill-posedness と次元の呪いが深刻になるため，$F$ による低次元構造を同時に利用する推定法が望ましい。

\subsection{Nonparametric regularization}

第二の実装方法は，d'Haultfoeuille (2010) と同様に，条件付き期待作用素を用いた逆問題として扱うことである。作用素 $T$ を
\begin{equation}
  T\varphi(v)=\E\{R\varphi(\Y)\mid V=v\}
\end{equation}
と定義すれば，$w$ は
\begin{equation}
  Tw=1
  \label{eq:operator-equation}
\end{equation}
を満たす。連続変数の場合，この逆問題は一般に ill-posed であるため，Tikhonov 正則化などを用いる。例えば，$\widehat T$ をカーネルまたは系列推定で得た作用素として，
\begin{equation}
  \widehat w\in
  \arg\min_{\varphi\in\mathcal W}
  \left\|\widehat T\varphi-1\right\|^2+\rho_N\|\varphi\|^2
  \label{eq:tikhonov-vector}
\end{equation}
を考える。ここで $\mathcal W$ は $\varphi\geq 1$ などの制約を含む関数空間であり，$\rho_N$ は正則化パラメータである。

\subsection{IPW estimation of target functionals}

$\widehat w$ が得られれば，任意の汎関数
\begin{equation}
  \theta_\varphi=\E\{\varphi(\Y,V)\}
\end{equation}
は
\begin{equation}
  \widehat\theta_\varphi=\frac{1}{N}\sum_{i=1}^N R_i\widehat w(\Y_i)\varphi(\Y_i,V_i)
  \label{eq:ipw-estimator}
\end{equation}
により推定できる。特に，$\varphi(\Y,V)=\Y$ とすれば平均ベクトル $\E(\Y)$ の推定量を得る。分布関数や密度を推定したい場合には，$\varphi$ を指示関数やカーネル関数に置けばよい。

\subsection{Latent-function estimation}

最後に，復元された $P(\Y,V)$ または IPW 重み付き標本を用いて，
\begin{equation}
  \Y=\bm h(F)+\eps,
  \qquad F=g(V)+U
\end{equation}
を推定する。実装としては，以下の二つが考えられる。
\begin{enumerate}[(i)]
  \item $g$ と $\bm h$ に series/spline 表現を置き，IPW 重み付き尤度または最小距離で推定する。
  \item $F$ を潜在変数として階層モデルを構成し，IPW 重み付き擬似尤度または Bayesian computation により推定する。
\end{enumerate}
本稿の主定理は識別結果であり，漸近理論は今後の課題とする。ただし，第一段階の $\widehat w$ の正則化誤差と，第二段階の潜在関数推定誤差を分離して評価する必要がある。

\section{Simulation studies}

本節では，理論結果を確認するためのシミュレーション設計を示す。ここでは雛形を記載し，具体的な設定は後続の実装で確定する。

\subsection{Data-generating process}

サンプルサイズを $N$ とし，$V_i$ を常時観測される shadow variable とする。例えば
\begin{align}
  V_i &\sim <V_distribution>,\\
  U_i &\sim <U_distribution>,\\
  F_i &= g(V_i)+U_i,\\
  Y_{ij} &= h_j(F_i)+\varepsilon_{ij},\qquad j=1,\ldots,K,\\
  R_i &=\1\{k(\Y_i)+\eta_i>0\}
\end{align}
によりデータを生成する。ここで $\Y_i=(Y_{i1},\ldots,Y_{iK})'$ である。欠測確率が $\Y_i$ 全体に依存するように $k(\cdot)$ を設定することで，$Y_j$ 間の相関と nonignorable missingness を同時に導入する。

\subsection{Compared methods}

比較対象は以下を想定する。
\begin{enumerate}[(i)]
  \item Full data: 欠測のない完全データに基づく推定。
  \item Complete case: $R=1$ の標本のみに基づく推定。
  \item MAR/IPW: 観測変数のみに基づき MAR を仮定した補正。
  \item Proposed method: $V$ を shadow variable として用いた提案法。
\end{enumerate}

\subsection{Evaluation metrics}

評価指標は，平均ベクトル $\E(\Y)$ のバイアス，MSE，および測定関数 $g,h_1,\ldots,h_K$ の推定誤差とする。例えば
\begin{equation}
  \mathrm{MSE}_{\Y}=\left\|\widehat{\E(\Y)}-\E(\Y)\right\|^2
\end{equation}
を用いる。関数推定については，格子点上の integrated squared error を用いる。

\section{Empirical illustration}

本節では，実証分析または擬似実証の雛形を示す。現時点ではデータ名や変数名が未確定であるため，未確定箇所はプレースホルダとして記載する。

\subsection{Data}

データセットとして $<data\_source\_name>$ を用いる。分析対象期間は $<sample\_period>$ である。ベクトル値アウトカムは
\begin{equation}
  \Y=(<data\_Y1\_name>,\ldots,<data\_YK\_name>)'
\end{equation}
であり，常時観測される shadow variable は
\begin{equation}
  V=<data\_V\_name>
\end{equation}
である。経済データを用いる場合，$V$ は景気先行指標，$\Y$ は複数企業または複数市場の指標として定義できる。

\subsection{Artificial nonignorable missingness}

完全データが利用可能な場合，まず擬似実証として人工的に NI 欠測を導入する。例えば
\begin{equation}
  R_i=\1\{\alpha_0+\alpha'\Y_i+\eta_i>0\}
\end{equation}
により，$\Y_i$ に依存する欠測を発生させる。欠測率は $<missing\_rate>$ となるように $\alpha_0$ を調整する。

\subsection{Estimation procedure}

提案法では，まず $\E\{Rw(\Y)\mid V\}=1$ に基づき $w(\Y)$ を推定する。次に，IPW により $P(\Y,V)$ および $\E(\Y)$ を推定する。最後に，復元された分布に基づき潜在測定関数 $g$ および $\bm h$ を推定する。

\subsection{Results}

結果は表 \ref{tab:empirical-template} の形式で示す予定である。

\begin{table}[tbp]
\centering
\caption{実証分析または擬似実証の結果雛形}
\label{tab:empirical-template}
\begin{tabular}{lccc}
\hline
Method & Bias of $\E(\Y)$ & MSE of $\E(\Y)$ & Notes\\
\hline
Full data & <value> & <value> & Benchmark\\
Complete case & <value> & <value> & Biased under NI\\
MAR/IPW & <value> & <value> & Uses observed covariates only\\
Proposed & <value> & <value> & Uses shadow variable $V$\\
\hline
\end{tabular}
\end{table}

\section{Conclusion}

本稿の結論は次の通りである。第一に，欠測がベクトル値の測定変数 $\Y$ に依存する
\begin{equation}
  R^*=k(\Y)+\eta
\end{equation}
の場合，$V$ が常に観測され，$R\indep V\mid \Y$ と vector completeness が成り立てば，$\pi(\Y)$，$P(\Y,V)$，および $\E(\Y)$ は識別される。

第二に，復元された joint distribution $P(\Y,V)$ に非線形測定誤差モデルの識別定理を適用すれば，IV型構造 $F=g(V)+U$ の下で $g$ と $h_1,\ldots,h_K$ が識別される。

第三に，旧稿の $Y_j$ ごとの識別ではなく，$\Y$ 全体に条件づける表記へ変更することで，$Y_j$ 間の相関と欠測指標間の依存を許した理論に整理できる。この点が，今回の会議で確認された主要な修正点である。

第四に，proxy 型 $V=g(F)+U$ は測定誤差モデルとして興味深いが，欠測補完の主定理としては識別の見通しがまだ明確でないため，本稿では中心的主張から外す。ただし，補足的な識別議論は付録にまとめる。

\appendix

\section{Proofs}

\subsection{Proof of Proposition \ref{prop:shadow-propensity-id}}
$w$ と $\widetilde w$ がともに \eqref{eq:shadow-integral-equation} を満たすとする。すると
\begin{align}
0
&=\E\{R[w(Y)-\widetilde w(Y)]\mid Z\}\\
&=\E\{\pi(Y)[w(Y)-\widetilde w(Y)]\mid Z\}.
\end{align}
仮定 \ref{ass:shadow-completeness-review} より $\pi(Y)[w(Y)-\widetilde w(Y)]=0$ a.s. である。positivity より $\pi(Y)>0$ なので $w(Y)=\widetilde w(Y)$ a.s. が従う。よって \eqref{eq:shadow-integral-equation} の解は一意であり，$w$ と $\pi$ は識別される。

\subsection{Proof of Lemma \ref{lem:vector-shadow-independence}}
$R$ は $\Y$ と $\eta$ のみの関数であり，$\eta$ は $V$ および他の誤差項から独立である。したがって
\begin{align}
  \Prb(R=1\mid \Y,V)
  &=\Prb\{\eta>-k(\Y)\mid \Y,V\}\\
  &=\Prb\{\eta>-k(\Y)\mid \Y\}\\
  &=\pi(\Y).
\end{align}
よって $R\indep V\mid \Y$ が従う。

\subsection{Proof of Theorem \ref{thm:vector-shadow-identification}}
$w(\bm y)=1/\pi(\bm y)$ と置く。補題 \ref{lem:vector-shadow-independence} より，
\begin{align}
  \E\{Rw(\Y)\mid V\}
  &=\E\left[\E\{Rw(\Y)\mid \Y,V\}\mid V\right]\\
  &=\E\left[w(\Y)\pi(\Y)\mid V\right]\\
  &=1.
  \label{eq:appendix-vector-ipw-moment}
\end{align}
左辺は $R=1$ のときのみ $\Y$ を用いるため，観測分布から評価できる。したがって $w$ は観測可能な積分方程式
\begin{equation}
  \E\{Rw(\Y)\mid V\}=1
\end{equation}
の解である。

次に一意性を示す。$w$ と $\widetilde w$ がこの積分方程式の二つの解であるとする。すると
\begin{align}
0
&=\E\{R[w(\Y)-\widetilde w(\Y)]\mid V\}\\
&=\E\{\pi(\Y)[w(\Y)-\widetilde w(\Y)]\mid V\}.
\end{align}
仮定 \ref{ass:vector-completeness} より $\pi(\Y)[w(\Y)-\widetilde w(\Y)]=0$ a.s. である。positivity より $w=\widetilde w$ a.s. が従う。よって $w$ および $\pi$ は識別される。

任意の有界可測関数 $\varphi$ について
\begin{equation}
  \E\{\varphi(\Y,V)\}
  =\E\{Rw(\Y)\varphi(\Y,V)\}
  \label{eq:appendix-vector-ipw-functional}
\end{equation}
が成り立つ。したがって $P(\Y,V)$ が識別される。特に $\varphi(\Y,V)=\Y$ とすれば平均ベクトル $\E(\Y)$ が識別される。

\subsection{Proof of Theorem \ref{thm:vector-iv-functions}}
定理 \ref{thm:vector-shadow-identification} より，$P(\Y,V)$ と $\E(\Y)$ は識別される。復元された $P(\Y,V)$ は
\begin{equation}
  p(\bm y\mid v)
  =\int
    \left\{\prod_{j=1}^K p_{\varepsilon_j}(y_j-h_j(f))\right\}
    p_U\{f-g(v)\}
    df
  \label{eq:appendix-vector-iv-kernel}
\end{equation}
という積分表現を持つ。仮定 \ref{ass:vector-latent-identification} により，この表現は正規化の下で一意である。したがって，復元された $P(\Y,V)$ から $g$ および $\bm h=(h_1,\ldots,h_K)$ が識別される。

\subsection{Proof of Corollary \ref{cor:two-dimensional}}
$K=2$ の場合，定理 \ref{thm:vector-shadow-identification} を $\Y=(Y_1,Y_2)'$ に適用すればよい。欠測確率は $\pi(Y_1,Y_2)=\Prb(R=1\mid Y_1,Y_2)$ であり，$\pi_1(Y_1)\pi_2(Y_2)$ と分解する必要はない。したがって，$Y_1$ と $Y_2$ の相関および欠測確率の同時依存を許したまま，$P(Y_1,Y_2,V)$ が識別される。

\section{Additional measurement error}
\label{app:additional-measurement-error}

観測変数が，真の測定値 $Y_j^*$ にさらに測定誤差 $Q_j$ を加えた
\begin{align}
  Y_j^*&=h_j(F)+e_j,\\
  Y_j&=Y_j^*+Q_j
\end{align}
であるとする。ここで $e_j,Q_j$ は互いに独立であり，$F$ や $V$ とも独立とする。

\begin{proposition}[誤差項の結合]
\label{prop:vector-error-combination}
$\varepsilon_j=e_j+Q_j$ と定義すれば，モデルは
\begin{equation}
  Y_j=h_j(F)+\varepsilon_j
\end{equation}
に帰着する。したがって，$\varepsilon_j$ が仮定 \ref{ass:vector-latent-identification} に必要な独立性と正則性を満たす限り，追加測定誤差は $\bm h$ と $g$ の識別可能性を変えない。
\end{proposition}

\begin{proof}
代入により
\begin{equation}
  Y_j=h_j(F)+e_j+Q_j=h_j(F)+\varepsilon_j
\end{equation}
である。独立な誤差の和は，$F,V$ および他の測定誤差から独立な新しい誤差項として扱える。よって，識別問題は基本モデル \eqref{eq:vector-measurement} に一致する。
\end{proof}

\section{Empirical Bayes interpretation and completeness}
\label{app:empirical-bayes-completeness}

誤差分布をノンパラメトリックに扱う場合でも，識別に必要なのは特定の分布族を仮定することではなく，誤差が情報を完全には消去しないことである。この点は completeness とフーリエ変換によって表現できる。

\begin{proposition}[シフト族における completeness の十分条件]
\label{prop:fourier-completeness}
$F=g(V)+U$ を考える。$U$ は密度 $p_U$ を持ち，$U\indep V$ とする。$\phi_U(t)=\E\{\exp(itU)\}$ を $U$ の特性関数とする。さらに $g(\calV)=\R$ とし，$\phi_U(t)\neq0$ が全ての $t\in\R$ で成り立つとする。このとき，
\begin{equation}
  \E\{a(F)\mid V\}=0\quad \mathrm{a.s.}
  \Longrightarrow
  a(F)=0\quad \mathrm{a.s.}
\end{equation}
が成り立つ。
\end{proposition}

\begin{proof}
$x=g(v)$ とおくと，$F\mid V=v$ の密度は $p_U(f-x)$ である。したがって
\begin{equation}
  0=\E\{a(F)\mid V=v\}=\int a(f)p_U(f-x)df=(a*\check p_U)(x),
\end{equation}
ただし $\check p_U(t)=p_U(-t)$ である。$g(\calV)=\R$ なので，$(a*\check p_U)(x)=0$ が全ての $x\in\R$ で成り立つ。フーリエ変換をとると
\begin{equation}
  \widehat a(t)\overline{\phi_U(t)}=0
\end{equation}
である。$\phi_U(t)\neq0$ より $\widehat a(t)=0$ であり，フーリエ変換の一意性から $a=0$ a.e. が従う。
\end{proof}

\begin{remark}[ベクトル $\Y$ に対する completeness]
本稿の主定理に必要なのは $\E\{a(\Y)\mid V\}=0\Rightarrow a(\Y)=0$ というベクトル版 completeness である。これはスカラー版より強い条件であり，実証では $V$ の次元，$g(V)$ の変動幅，および $\Y$ の次元に注意が必要である。理論上はスカラーからベクトルへの拡張は可能だが，仮定の強さと推定の ill-posedness は明示しておく必要がある。
\end{remark}

\begin{remark}[測定誤差と欠測の境界]
経験ベイズ的には，誤差分布を推定できるかどうかは，情報が単に拡散しているのか，あるいは完全に遮断されているのかの違いとして理解できる。測定誤差は畳み込みによる smoothing であり，特性関数が消えなければ deconvolution により復元可能である。一方，欠測や切断によりある領域の情報が完全に失われると，対応する積分作用素は非自明なヌル空間を持ち，completeness は破綻する。したがって，「誤差分布をノンパラメトリックに設定できること」は，「作用素が十分なランクを持つこと」と同じ方向の条件である。
\end{remark}

\section{Missingness depending on the latent factor}
\label{app:latent-factor-missingness}

次に，欠測が $\Y$ ではなく潜在因子 $F$ に依存する場合を考える：
\begin{equation}
  R=\1\{R^*>0\},
  \qquad
  R^*=k(F)+\eta.
  \label{eq:vector-f-missing}
\end{equation}
このとき $\rho(f)=\Prb(R=1\mid F=f)$ と書く。

\begin{proposition}[shadow 条件の破綻]
\label{prop:vector-f-missing-not-shadow}
欠測機構 \eqref{eq:vector-f-missing} の下では，一般に
\begin{equation}
  R\not\indep V\mid \Y.
\end{equation}
したがって，$V$ を $\Y$ に対する shadow variable とする定理 \ref{thm:vector-shadow-identification} の証明は適用できない。
\end{proposition}

\begin{proof}
\eqref{eq:vector-f-missing} より
\begin{equation}
  \Prb(R=1\mid \Y=\bm y,V=v)=\E\{\rho(F)\mid \Y=\bm y,V=v\}.
\end{equation}
$\Y$ は $F$ の誤差付き測定値であるため，$\Y=\bm y$ を条件づけても $F$ は一般に確定しない。さらに $V$ は $F$ と関連しているので，$F\mid \Y=\bm y,V=v$ の分布は一般に $v$ に依存する。$\rho$ が非定数ならば右辺も一般に $v$ に依存する。よって $R\indep V\mid \Y$ は成立しない。
\end{proof}

\begin{proposition}[母集団平均の非識別]
\label{prop:mean-nonid-f-missing}
欠測が \eqref{eq:vector-f-missing} に従う場合，追加のパラメトリック制約なしには，母集団平均 $\E(\Y)$ は一般にノンパラメトリック識別されない。
\end{proposition}

\begin{proof}
観測される完全ケース分布は，概念的には
\begin{equation}
  p_{\mathrm{obs}}(\bm y,v)
  \propto
  \int
    \left\{\prod_{j=1}^K p_{\varepsilon_j}(y_j-h_j(f))\right\}
    p_U\{f-g(v)\}
    \rho(f)p_F(f)df
\end{equation}
の形で表される。観測分布は $p_F(f)$ と $\rho(f)$ を別々には含まず，積 $\rho(f)p_F(f)$ を通じてのみ含む。任意の正の関数 $a(f)$ によって $p_F$ を傾け，$\rho$ を逆方向に調整すれば，観測分布を保ったまま母集団分布を変えることができる。したがって $\E(\Y)$ は一般に一意に定まらない。
\end{proof}

\section{Proxy-type approach}
\label{app:proxy-type}

本付録では，旧稿で扱っていた proxy 型構造
\begin{equation}
  V=g(F)+U
  \label{eq:app-proxy-structure}
\end{equation}
を補足的に整理する。本文の主定理は IV 型構造 $F=g(V)+U$ に集中する。したがって，以下の結果は本稿の中心的主張ではなく，追加の多重指標識別条件を置いた場合の参考結果である。

\subsection{Shadow variable condition under proxy-type structure}

proxy 型では，$V$ は $F$ の原因ではなく，$F$ の誤差付き測定値である。すなわち，
\begin{equation}
  \Y=\bm h(F)+\eps,
  \qquad
  V=g(F)+U
  \label{eq:app-proxy-measurement}
\end{equation}
を考える。ここで $F,U,\eps,\eta$ は互いに独立であり，欠測機構は
\begin{equation}
  R=\1\{R^*>0\},
  \qquad
  R^*=k(\Y)+\eta
  \label{eq:app-proxy-y-missing}
\end{equation}
に従うとする。

このとき，$R$ は $\Y$ と $\eta$ のみの関数であり，$\eta$ は $V$ から独立であるため，
\begin{equation}
  R\indep V\mid \Y
  \label{eq:app-proxy-shadow}
\end{equation}
が成立する。したがって，positivity と vector shadow completeness
\begin{equation}
  \E\{a(\Y)\mid V\}=0
  \quad\Longrightarrow\quad
  a(\Y)=0
  \label{eq:app-proxy-completeness}
\end{equation}
が成り立つならば，本文の定理 \ref{thm:vector-shadow-identification} と同じ議論により，$\pi(\Y)$ と $P(\Y,V)$ は識別される。

\subsection{Supplementary identification result}

さらに，復元された $P(\Y,V)$ に対して，多重指標モデル
\begin{equation}
  p(\bm y,v)
  =
  \int
  \left\{\prod_{j=1}^K p_{\varepsilon_j}(y_j-h_j(f))\right\}
  p_U\{v-g(f)\}
  p_F(f)df
  \label{eq:app-proxy-kernel}
\end{equation}
が，正規化の下で一意に識別されると仮定する。この条件は，Kotlarski 型の補題や非線形多重指標モデルの識別定理に対応する。少なくとも三つ以上の独立指標が存在し，位置・尺度・符号の正規化が与えられることが典型的に必要である。

\begin{proposition}[proxy 型における補足的識別結果]
\label{prop:app-proxy-identification}
\eqref{eq:app-proxy-measurement}，\eqref{eq:app-proxy-y-missing}，positivity，vector shadow completeness，および多重指標モデル \eqref{eq:app-proxy-kernel} の識別条件が成り立つとする。このとき，$P(\Y,V)$，$\E(\Y)$，$g$，および $h_1,\ldots,h_K$ は観測分布から識別される。
\end{proposition}

\begin{proof}
まず，\eqref{eq:app-proxy-shadow}，positivity，vector shadow completeness により，本文の定理 \ref{thm:vector-shadow-identification} と同じ積分方程式
\begin{equation}
  \E\{Rw(\Y)\mid V\}=1,
  \qquad
  w(\Y)=1/\pi(\Y)
\end{equation}
の解が一意に定まる。したがって，$\pi(\Y)$，$P(\Y,V)$，および $\E(\Y)$ は識別される。

次に，復元された $P(\Y,V)$ は \eqref{eq:app-proxy-kernel} の積分表現を持つ。多重指標モデルの識別条件により，この表現は正規化の下で一意である。よって，$g$ と $h_1,\ldots,h_K$ が識別される。
\end{proof}

\begin{remark}[本文から付録へ移した理由]
proxy 型では $V$ が $F$ の追加測定値として機能するため，測定誤差モデルとしては自然である。一方で，本稿の主な貢献である「常時観測補助変数によりベクトル値欠測変数を補完する」という解釈は，IV 型構造 $F=g(V)+U$ の方が明確である。そのため，proxy 型の識別は本稿では補足的議論として付録に留める。
\end{remark}

\nocite{dHaultfoeuille2010,HuSchennach2008,Kotlarski1967,ZhaoShao2015,KanoTakai2011,WangShaoKim2014,ShaoWang2016,MiaoTchetgen2016,MiaoTchetgen2018,LiMiaoTchetgen2023,Rubin1976,Rubin1987,LittleRubin2002,BlackwellHonakerKing2017a,BlackwellHonakerKing2017b}
\printbibliography[title={References}]

\iffalse
\begin{thebibliography}{99}

\bibitem{dHaultfoeuille2010}
D'Haultfoeuille, X. (2010).
A new instrumental method for dealing with endogenous selection.
\textit{Journal of Econometrics}, 154(1), 1--15.

\bibitem{HuSchennach2008}
Hu, Y. and Schennach, S. M. (2008).
Instrumental variable treatment of nonclassical measurement error models.
\textit{Econometrica}, 76(1), 195--216.

\bibitem{Kotlarski1967}
Kotlarski, I. I. (1967).
On characterizing the gamma and the normal distribution.
\textit{Pacific Journal of Mathematics}, 20(1), 69--76.

\bibitem{ZhaoShao2015}
Zhao, J. and Shao, J. (2015).
Semiparametric pseudo-likelihoods in generalized linear models with nonignorable missing data.
\textit{Journal of the American Statistical Association}, 110(512), 1577--1590.

\bibitem{WangShaoKim2014}
Wang, S., Shao, J. and Kim, J. K. (2014).
An instrumental variable approach for identification and estimation with nonignorable nonresponse.
\textit{Statistica Sinica}, 24, 1097--1116.

\bibitem{ShaoWang2016}
Shao, J. and Wang, L. (2016).
Semiparametric inverse propensity weighting for nonignorable missing data.
\textit{Biometrika}, 103(1), 175--187.

\bibitem{MiaoTchetgen2016}
Miao, W. and Tchetgen Tchetgen, E. J. (2016).
On varieties of doubly robust estimators under missingness not at random with a shadow variable.
\textit{Biometrika}, 103(2), 475--482.

\bibitem{MiaoTchetgen2018}
Miao, W. and Tchetgen Tchetgen, E. J. (2018).
Identification and inference with nonignorable missing covariate data.
\textit{Statistica Sinica}, 28(4), 2049--2067.

\bibitem{LiMiaoTchetgen2023}
Li, W., Miao, W. and Tchetgen Tchetgen, E. J. (2023).
Non-parametric inference about mean functionals of non-ignorable non-response data without identifying the joint distribution.
\textit{Journal of the Royal Statistical Society: Series B}, 85(3), 913--935.

\bibitem{Rubin1976}
Rubin, D. B. (1976).
Inference and missing data.
\textit{Biometrika}, 63(3), 581--592.

\bibitem{Rubin1987}
Rubin, D. B. (1987).
\textit{Multiple Imputation for Nonresponse in Surveys}.
Wiley.

\bibitem{LittleRubin2002}
Little, R. J. A. and Rubin, D. B. (2002).
\textit{Statistical Analysis with Missing Data}.
Wiley.

\bibitem{BlackwellHonakerKing2017a}
Blackwell, M., Honaker, J. and King, G. (2017a).
A unified approach to measurement error and missing data: Overview and applications.
\textit{Sociological Methods \& Research}, 46(3), 303--341.

\bibitem{BlackwellHonakerKing2017b}
Blackwell, M., Honaker, J. and King, G. (2017b).
A unified approach to measurement error and missing data: Details and extensions.
\textit{Sociological Methods \& Research}, 46(3), 342--369.

\end{thebibliography}
\fi

\end{document}
